\subsection{Proof by Exhaustion}
\label{subsections:techniques:universal:exhaustion}



\paragraph{Description.}
A proof by exhaustion proves a mathematical statement by checking every %
possible case. %
Precisely, it proves a universal statement, %
$\forall s \in S,\ P(s)$, where $S$ is finite, by listing every element %
of $S$ and showing that $P(s)$ is true for each one. %
It is effectively a %
\hyperref[subsections:techniques:universal:cases]{%
    proof by division into cases}, %
but with specific elements instead of broad definitions for elements. %
Generally, a proof by exhaustion mirrors the following format:
\begin{enumerate}{
    \item \textbf{Setup.} %
        Define $S = \{s_1, \dots, s_n\}$ and let $s \in S$.
    \item \textbf{Verification.} %
        For each $s_i$, suppose $s = s_i$ and show that $P(s_i)$ is true.
    \item \textbf{Conclusion.} %
        Conclude the proof, optionally by summarizing what has been found.
}
\end{enumerate}



\paragraph{Advice.}
\begin{itemize}{
    \item Use this proof only for a universal statement, %
        $\forall s \in S,\ P(s)$, when $S$ is finite and small enough to %
        check element by element. %
        This proof cannot be used when $S$ is infinite.
    \item Show $P(s_i)$ for every element of $S$, as any unaccounted for %
        element leaves the statement unproven.
}
\end{itemize}



\paragraph{Example.}
\begin{exercise}{
    Prove that $\forall k \in \{1, 2\},\ k^2 < 5$.
}
\end{exercise}
\begin{proof}{
    This statement can be proven by exhaustion.

    \begin{annotatedpanel}[Setup.]{
        Let $K = \{1, 2\}$. %
        Let $k \in K$.
    }{
        \textbf{Step 1.} Define $S$ and let $s \in S$. %
        Here, $S = K = \{1, 2\}$, so $s_1 = 1$ and $s_2 = 2$, %
        and $s = k$.
    }
    \end{annotatedpanel}

    \begin{annotatedpanel}[Verification.]{
        Suppose $k = 1$. %
        Then $k^2 = 1^2 = 1 < 5$. \\
        Suppose $k = 2$. %
        Then $k^2 = 2^2 = 4 < 5$.
    }{
        \textbf{Step 2.} For each $s_i$, suppose $s = s_i$ and show that %
        $P(s_i)$ is true. %
        Here, $P(k)$ is the inequality $k^2 < 5$.
    }
    \end{annotatedpanel}

    \begin{annotatedpanel}[Conclusion.]{
        Thus, $\forall k \in K,\ k^2 < 5$. %
        Therefore, the statement holds by exhaustion. \qedhere
    }{
        \textbf{Step 3.} Conclude the proof. %
        Here, the summary states what has been found: %
        $P(k)$ holds for both elements.
    }
    \end{annotatedpanel}
}
\end{proof}